%-----------------------------------LICENSE------------------------------------%
%   This file is part of Mathematics-and-Physics.                              %
%                                                                              %
%   Mathematics-and-Physics is free software: you can redistribute it and/or   %
%   modify it under the terms of the GNU General Public License as             %
%   published by the Free Software Foundation, either version 3 of the         %
%   License, or (at your option) any later version.                            %
%                                                                              %
%   Mathematics-and-Physics is distributed in the hope that it will be useful, %
%   but WITHOUT ANY WARRANTY; without even the implied warranty of             %
%   MERCHANTABILITY or FITNESS FOR A PARTICULAR PURPOSE.  See the              %
%   GNU General Public License for more details.                               %
%                                                                              %
%   You should have received a copy of the GNU General Public License along    %
%   with Mathematics-and-Physics.  If not, see <https://www.gnu.org/licenses/>.%
%------------------------------------------------------------------------------%
\documentclass{article}
\usepackage{amssymb}
\usepackage{amsmath}

\title{18.100P: Final}
\date{Spring 2026}

% No indent and no paragraph skip.
\setlength{\parindent}{0em}
\setlength{\parskip}{0em}

\begin{document}
    \maketitle
    \textbf{Series / Uniform Continuity} (10 pts)
    \par
    Suppose $a:\mathbb{N}\rightarrow\mathbb{R}$ is a sequence of
    non-negative real numbers ($a_{n}\geq{0}$ for all $n\in\mathbb{N}$),
    and suppose we have a sequence of
    functions $F_{n}:\mathbb{R}\rightarrow\mathbb{R}$ such that for
    all $x\in\mathbb{R}$ and for all $n\in\mathbb{N}$
    we have $|F_{n}(x)|\leq{a}_{n}$, and suppose
    $S_{N}$ denotes the $N^{\textrm{th}}$ partial sum:

    \begin{equation}
        S_{N}(x)
        =
        \sum_{n=0}^{N}
            F_{n}(x).
    \end{equation}

    Prove that if

    \begin{equation}
        \lim_{N\rightarrow\infty}
        \sum_{n=0}^{N}
            a_{n}
        <
        \infty,
    \end{equation}

    then there is a function $f:\mathbb{R}\rightarrow\mathbb{R}$
    such that $S_{N}(x)\rightarrow{f}(x)$ uniformly.
    That is, if the sum of the $a_{n}$ converges, then the sequence
    of partial sums $S_{N}$ convergences uniformly.
    \newpage
    \textbf{Differentiation}
    \begin{itemize}
        \item (5 pts)
            Differentiability is a point-wise notion.
            Prove that

            \begin{equation}
                f(x)
                =
                \begin{cases}
                    x^{2},&x\in\mathbb{Q}\\
                    -x^{2},&x\not\in\mathbb{Q}
                \end{cases}
            \end{equation}

            is differentiable at $x=0$ and nowhere else.
        \item (5 pts)
            Suppose $g:\mathbb{R}\rightarrow\mathbb{R}$ is differentiable
            everywhere. Prove that $g$ is continuous everywhere.
    \end{itemize}
    \newpage
    \textbf{Differentiation in Higher Dimensions} (10 pts)
    \par
    Differentiation can be extended to higher dimensions.
    The \textbf{total derivative} of a function
    $f:\mathcal{U}\rightarrow\mathbb{R}^{m}$, defined on an open set
    $\mathcal{U}\subseteq\mathbb{R}^{n}$, at a point
    $\mathbf{x}_{0}\in\mathcal{U}$ is the unique linear function
    $%
        \operatorname{d}f_{\mathbf{x}_{0}}:%
        \mathbb{R}^{n}\rightarrow\mathbb{R}^{m}%
    $
    (if it exists) such that:

    \begin{equation}
        \lim_{\mathbf{x}\rightarrow\mathbf{x}_{0}}
        \frac{%
            \|f(\mathbf{x})-f(\mathbf{x}_{0})-
                \operatorname{d}f_{\mathbf{x}_{0}}
                \left(\mathbf{x}-\mathbf{x}_{0}\right)
            \|%
        }{%
            \|\mathbf{x}-\mathbf{x}_{0}\|%
        }=0
    \end{equation}

    A quadratic form $q:\mathbb{R}^{n}\rightarrow\mathbb{R}$
    is a function such that there is an $n\times{n}$ matrix $A$
    with $q(\mathbf{x})=\mathbf{x}^{T}A\mathbf{x}$ where
    $\mathbf{x}^{T}$ is the transpose of $\mathbf{x}$.
    \par\hfill\par
    Compute the total derivative of $q$
    at a point $\mathbf{x}_{0}\in\mathbb{R}^{n}$.
    \par\hfill\par
    Hint: The following facts from linear algebra may help:

    \begin{equation}
        \begin{array}{rcl}
            \displaystyle
            (A+B)^{T}
            &=&
            \displaystyle
            A^{T}+B^{T}\\[1.5em]
            (AB)^{T}
            &=&
            B^{T}A^{T}.
        \end{array}
    \end{equation}

    Since $\textrm{d}\,q_{\mathbf{x}_{0}}$ is linear, it is given by a
    matrix.
    \newpage
    \textbf{Riemann Integration} (10 pts)
    \par
    Suppose $f:[a,\,b]\rightarrow\mathbb{R}$ is monotonically
    increasing, and that it is bounded above and below.
    $f$ does \textbf{not} need to be continuous. Prove $f$
    is Riemann / Darboux integrable.
    \newpage
    \textbf{Fourier Series} (10 pts)
    \par
    Define $f:[0,\,1]\rightarrow\mathbb{R}$ via $f(x)=x$. Extending
    $f$ periodically to $\mathbb{R}$ produces a
    \textit{sawtooth wave}. Compute the coefficients for the
    Fourier sine series of $f$.
    \newpage
    \textbf{Fourier Analysis}
    \begin{itemize}
        \item (5 pts)
            Prove or disprove the following.
            There is a Riemann integrable function
            $f:[-1,\,1]\rightarrow\mathbb{R}$ such that the
            Fourier series of $f$ is given by

            \begin{equation}
                f(x)
                =
                \sum_{n=0}^{\infty}
                    (-1)^{n}
                    \sin(n\pi{x}).
            \end{equation}

        \item (5 pts)
            Prove or disprove the following.
            There is a Riemann integrable function
            $f:[-1,\,1]\rightarrow\mathbb{R}$ such that the
            Fourier series of $f$ is given by

            \begin{equation}
                f(x)
                =
                \sum_{n=0}^{\infty}
                    \frac{1}{1+n^{2}}
                    \cos(n\pi{x}).
            \end{equation}
    \end{itemize}
\end{document}
